%!TEX root = ../thesis.tex

\chapter{Related Work}\label{chap:related-work}

Fuzzing is a sequential decision making problem under certainty over a large search space
with a scalar progress signal.
This makes it amenable to solving it through reinforcement learning methods.
We will introduce and classify selected prior works in this area.

\textcite{wang_systematic_2020} classify ML-based fuzzers into five categories according
to the problems they solve.
In addition, it divides each fuzzer into what learning algorithm it uses.
These are broadly divided into traditional machine learning, deep learning, and reinforcement learning.
Reinforcement learning-based fuzzers they survey for each fuzzing step are shown in~\cref{tab:rw-rl-fuzzers}.

%!TEX root = ../thesis.tex

% Source: Wang et al., "A systematic review of fuzzing based on machine learning
% techniques", PLOS ONE 15(8), 2020, Table 2 (primary studies) and Table 3 (studies
% per fuzzing step). Rows below are the primary studies whose algorithm column names
% a reinforcement-learning method; the survey's study identifiers are kept here for
% traceability: SP1 Becker, SP13 Karamcheti, SP17 Boettinger, SP19 Drozd,
% SP22 Fang, SP30 Liu, SP35 Kuznetsov.
%
% TODO: two rows still lack a bib entry. Add them to Zotero, then put the \cite
% next to the author-year label as done for the others:
%   SP13 Karamcheti, Mann and Rosenberg, "Adaptive Grey-Box Fuzz-Testing with
%        Thompson Sampling", AISec 2018 (survey ref. [52]).
%   SP22 Fang and Yan, "Emulation-Instrumented Fuzz Testing of 4G/LTE Android
%        Mobile Devices Guided by Reinforcement Learning", ESORICS 2018
%        (survey ref. [61]). The fuzzer is called LEFT.

\begin{table}[htb]
  \centering
  \caption[\gls{rl}-based fuzzers by fuzzing step]{The reinforcement-learning-based
    fuzzers surveyed by \textcite{wang_systematic_2020}, arranged by the fuzzing
  step they address.}
  \label{tab:rw-rl-fuzzers}
  \begin{tabular}{@{}lll@{}}
    \toprule
    Study & Fuzzer type & Algorithm \\
    \midrule
    \multicolumn{3}{@{}l@{}}{\itshape Mutation operator selection} \\
    Becker (2010)~\cite{hutchison_autonomic_2010}    & network  & \gls{sarsa} \\
    Karamcheti (2018)~\cite{karamchetiAdaptiveGreyBoxFuzzTesting2018} & software & Thompson sampling \\
    Böttinger (2018)~\cite{bottinger_deep_2018}      & software & Deep Q-learning \\
    Drozd (2018)~\cite{drozd_fuzzergym_2018}         & software & Deep double Q-learning \\
    Fang (2018)~\cite{lopez_emulation-instrumented_2018}                                      & network  & Q-learning \\
    Liu (2019)~\cite{liu_reinforcement_2019}         & software & Deep Q-learning \\
    \addlinespace
    \multicolumn{3}{@{}l@{}}{\itshape Exploitability analysis} \\
    Kuznetsov (2019)~\cite{kuznetsov_automated_2019} & software & Deep Q-learning \\
    \addlinespace
    \multicolumn{3}{@{}l@{}}{\itshape Seed file generation, testcase generation,
      testcase filtering} \\
    \multicolumn{3}{@{}l@{}}{No reinforcement-learning-based work surveyed.} \\
    \bottomrule
  \end{tabular}
\end{table}


\section{Modelling fuzzing as a Markov decision process}\label{sec:rw-mdp-formulations}

For the six of the seven papers that use a form of Q-Learning or SARSA, we
outline how they model their corresponding problem as an \gls{mdp} and
summarize the state, action and reward design.

\begin{description}[style=nextline, leftmargin=0pt, labelindent=0pt, labelwidth=0pt]
  \item[Becker et al.\ (2010)~\cite{hutchison_autonomic_2010}]
    They fuzz the IPv6 Neighbor Discovery Protocol with \gls{sarsa}.

    \begin{itemize}
      \item \textbf{State.}
        States are modeled as six different fuzzing strategies, thereby
        avoiding the state explosion problem.
        Three strategies mutate the content of a message (delete a field, insert a
        field, modify a field value) and three mutate the message order within a
        session (insert, repeat, or drop a message according to the behavioural
        model).

      \item \textbf{Action.}
        Actions are transitions to other fuzzing strategies.

      \item \textbf{Reward.}
        The reward consists of a sum of three different rewards
        \[
          r(q_t) = r_{\text{trace}}(q_t) + r_{\text{deb}}(q_t) + r_{\text{mon}}(q_t),
        \]

        where $r_{\text{trace}}$  depends on the entropy of a message and
        amount of function calls, $r_{\text{deb}}$ depends on whether
        an error was seen, and $r_{\text{mon}}$ depends on whether a corrupted
        or delayed messages was monitored.
    \end{itemize}
    The paper only provides a design and no working proof of concept and therefore
    no evaluation and empirical studies.

  \item[Böttinger et al.\ (2018)~\cite{bottinger_deep_2018}]
    They provided the first formulation of fuzzing as an \gls{mdp} and applied \gls{dqn}
    to fuzzing various PDF parsers.

    \begin{itemize}
      \item \textbf{State.}
        With $\Sigma$ a finite alphabet, all possible program inputs are given
        by the Kleene closure of $\Sigma$.
        They define the state space to be equal to the program input space $\mathcal{S}=\mathcal{I}=\Sigma^*$.

      \item \textbf{Action.}
        Actions are random variables mapping inputs to probabilistic rewriting rules.
        The implemented rules contain random bit flips, dictionary token insertion,
        input shifting, shuffling, and copying.

      \item \textbf{Reward.}
        The reward consists of a sum of two different rewards
        \[
          r(x, a) = E(x) + G(a),
        \]

        where $E$ depends on runtime properties of an execution trace $x$ $G$
        depends on the chosen action $a$. For $E$ they evaluate the number of
        newly discovered \glspl{bb}, the execution time, and the sum of both.
    \end{itemize}
    They found that neither experience replay nor a persistent history of discovered \glspl{bb}
    improve on a random baseline.

  \item[Drozd and Wagner (2018)~\cite{drozd_fuzzergym_2018}]
    FuzzerGym combines libFuzzer with OpenAI Gym, trained with deep double Q-learning and
    prioritised replay (\cref{fig:rw-fuzzergym}).

    \begin{itemize}
      \item \textbf{State.}
        They model the whole test input with a maximum input length of
        \num{4096} bytes as the state space, encoded as a bit-array
        $\mathcal{S}=\{0,1\}^{32768}$.

      \item \textbf{Action.}
        libFuzzer's thirteen mutation operators are selected as the action space.

      \item \textbf{Reward.}
        Reward is modeled as newly unique line coverage:
        \[
          R_t = \mathit{cov}_t - \mathit{cov}_{t-1},
        \]
        where $\mathit{cov}_t$ denotes the number of unique lines covered at time $t$.

    \end{itemize}
    To maintain the execution rate of libFuzzer of around \num{100000} executions per
    second, the authors let the fuzzer loop over a fixed-size circular buffer
    while new actions are only chosen every n'th fuzzing iteration, mirroring frame
    skipping in Atari environments.
    The problem the agent solves is therefore a \gls{pomdp} for which the authors
    employ an LSTM.

  \item[Fang and Yan (2018)~\cite{lopez_emulation-instrumented_2018}]
    LEFT (LTE-Oriented Emulation-Instrumented Fuzzing Testbed) is a fuzzer for the
    4G/LTE protocol stack on Android devices uilizing a network emulator.

    \begin{itemize}
      \item \textbf{State.}
        States are modeled as message types which the protocol implementation switches on.

      \item \textbf{Action.}
        An action is the switch branch the protocol dispatcher takes on an incoming message.

      \item \textbf{Reward.}
        The reward is derived from three sensor-driven signals:

        \[
          \beta^{\mathit{leak}}_c = 100, \qquad
          \beta^{\mathit{loss}}_a = 100, \qquad
          \beta_{\mathit{crash}} = -100,
        \]

        received when a protocol layer observes undecrypted data (\emph{leak}),
        when the device can no longer attach to a separate network after the episode (\emph{loss}),
        and when the emulator crashes (\emph{crash}).
    \end{itemize}

  \item[Liu et al.\ (2019)~\cite{authors_reinforcement_nodate}]
    They introduce FuzzBoost, a compiler fuzzer and automatic code synthesis framework.

    \begin{itemize}
      \item \textbf{State.}
        Similar to~\textcite{bottinger_deep_2018}, states are program inputs
        $\mathcal{S}=\mathcal{I}=\Sigma^{*}$.

      \item \textbf{Action.}
        Actions are also defined as mappings from program inputs to probabilistic rewriting rules,
        but are split into two families.
        Content rewrites operate at a token level and are aware of the grammar of the
        targeted programming language to overcome frontend checks.
        Window rewrites shift the input and grow or shrink it.
        A third, explicit \emph{terminate} action lets the agent end the episode itself.

      \item \textbf{Reward.}
        The reward for a given synthesized program $p$ is defined in relation to
        the entire fuzzing configuration $I$ as the ratio of unique basic blocks
        covered by $p$ to unique basic blocks covered by $I$:

        \[
          R(p,I)=\frac{B(T_p)}{\cup_{\rho\in I} B(T_p)}
        \]
    \end{itemize}
    They use an episodic framing instead of a continuous one like Böttinger.

  \item[Kuznetsov et al.\ (2019)~\cite{kuznetsov_automated_2019}]
    They provide a general architecture for applying deep RL to automated vulnerability
    testing.

    \begin{itemize}
      \item \textbf{State.}
        The authors don't describe a specific state model, but stay general and
        model \enquote{received data} as the state.
        This includes \enquote{program execution time, code coverage, code completion, etc.}.

      \item \textbf{Action.}
        Actions are mutations acting on a single seed.
        They are kept in a dictionary of size \num{45} and include operations such as
        lengthening and shortening strings, inserting integers or appending format specifiers.

      \item \textbf{Reward.}
        A reward is granted when execution time of the mutated seed exceeds the
        execution time of the previous seed or if a crash was encountered.

    \end{itemize}
    They evaluate their approach on synthetic targets.

\end{description}

Most of the surveyed papers use a similar approach.
They combine some form of Q-Learning with the program input as a state model
and an arbitrary list of mutations as actions.
Whenever the state is not modeled by program inputs, it is instead modeled by
handcrafted features.
This thesis aims to address these issues as follows:

\begin{itemize}
  \item \textbf{Better deep RL algorithms.}
    While the line of DQN has incorporated many improvements along the
    way~\cite{hessel_rainbow_2017}, the most impressive results have been
    achieved by some kind of MuZero~\cite{schrittwieser_mastering_2020}-style
    architecture. We propose using EfficientZero~\cite{wang_efficientzero_2024}
    as the core RL algorithm for improved learning and sampling efficiency.

  \item \textbf{Coverage as the state model.}
    \Gls{cgf} makes use of program coverage mostly as a way to decide whether to keep
    or discard test cases.
    This throws away precious information.
    We propose to include coverage information into the observations that the RL agent receives.
    We expect this information to aid the agent in decision making and therefore improve fuzzing performance.

  \item \textbf{Generic mutations.}
    Instead of choosing from a list of handcrafted mutations, we allow the agent
    to specify a location to mutate and a concrete value to set the location to.
    In theory, this allows the agent to learn the best mutation operators applicable
    to that specific PUT.
    Whether this beats a curated list of fixed mutation operators is an open question
    we plan to investigate empirically.

  \item \textbf{Generic reward signal.}
    A big advantage of RL-based approaches is that they don't rely on binary decisions,
    but can flexibly shape the reward signal in a way that is conducive to the fuzzing goal.
    One could easily imagine a reward signal that includes promoting the execution of
    specific system calls or regions of code.

\end{itemize}


Three patterns run through the table.
The state is almost always the input --- a substring window
\cite{bottinger_deep_2018, authors_reinforcement_nodate} or the whole test case as a
bit array \cite{drozd_fuzzergym_2018} --- and where it is not, it is a hand-built
abstraction: of the fuzzing strategy \cite{hutchison_autonomic_2010}, of the
protocol's message alphabet \cite{lopez_emulation-instrumented_2018}, or of the
program's response \cite{kuznetsov_automated_2019}.
The actions are, with hardly any exception, the mutation operators the underlying
fuzzer already had --- in LEFT, the message handlers the emulator already had ---
and the learning replaces the uniform random choice among them and nothing else.
And the reward is coverage, occasionally traded off against execution time
\cite{bottinger_deep_2018, kuznetsov_automated_2019} or against crash and
error indicators \cite{hutchison_autonomic_2010}; only LEFT dispenses with a
coverage term entirely, its target being a black box that cannot be instrumented.

The gap this leaves is the one this thesis addresses.
None of these agents observes the state of the \gls{put} itself, so none of them can
represent where in a stateful protocol or a multi-message session the target
currently is --- Becker et al.\ come closest, and explicitly discard the protocol
automaton to avoid a state explosion they have no machinery to handle.
LEFT gets nearest to a protocol state and still falls short of one, for the same
reason: a message type is not a position in the automaton, and the handset whose
state decides whether a perturbation lands is never observed at all.
None of them learns a model of the environment; every one is model-free, and pays for
it in sample efficiency, which is precisely what makes the millisecond-scale forward
pass of \textcite{drozd_fuzzergym_2018} such a problem that the \gls{mdp} has to be
degraded to a \gls{pomdp} to hide it.

% NOTE: the survey's summary of two of these formulations does not match the papers.
% It reports Becker's state as the protocol message type (the paper uses the fuzzing
% strategies as states, precisely to avoid the protocol automaton), and it groups
% Kuznetsov with the input-substring states (the paper's state is the program's
% response). The descriptions above follow the primary sources.

